\chapter{Outputs and the study report}
\label{ch:outputs}

\section{Files written by the trainer}

For a database \code{sqlite:///NAME.db} the trainer writes under \file{analysis/NAME/}.

\begin{description}
\item[\file{NAME.db}]
The Optuna study, present only when a search ran.
A fixed-architecture run does not create it.
SQLite may also leave \file{NAME.db-wal} and \file{NAME.db-shm} while a process has the file open.
A \file{NAME.db.lock} file serializes creation of a new database.

\item[\file{data/best_hyperparameters.txt}]
Family, width, depth, learning rate, PDE weights, and, for a quantum family, qubits, layers, and embedding.
Then the total, classical, and quantum trainable counts.
If a search ran, the Optuna trial number, its objective, and the stored loss pieces.
If L-BFGS-B ran, its success flag, message, iteration count, and evaluation count.
A skipped L-BFGS-B stage is recorded as skipped.

\item[\file{data/best_optuna_model.weights.h5}]
Keras weights of the final model.

\item[\file{data/best_optuna_model_config.json}]
The \code{HyperParams} record: width, depth, learning rate, both PDE weights, both log-weights, family, qubits, layers, and embedding.

\item[\file{data/<problem>_losses_data.npz}]
Adam total, mass, and momentum histories, and the same three series for L-BFGS-B when that stage ran.
The metadata repeats the family and the L-BFGS status.
\end{description}

Field archives depend on the benchmark.
One-dimensional problems write

\begin{itemize}
\item \file{<problem>_elevation_compare_data.npz} with arrays \code{TIME}, \code{X}, \code{A} (reference), and \code{M} (model);
\item \file{<problem>_velocity_compare_data.npz} with the same layout for \(u\);
\item \file{<problem>_error_data.npz} with \code{ETA_ERR} and \code{U_ERR};
\item \file{<problem>_profiles_boundary_data.npz} with elevation profiles and the velocity at the left endpoint.
\end{itemize}

The tide archives are converted to dimensional elevation, velocity, kilometers, and hours.
The dam break and the bump stay in the units of their dataclasses.
Steady square problems write \file{elevation}, \file{velocity_u}, and \file{velocity_v} compare files, plus \file{<problem>_error_data.npz} with \code{H_ERR}, \code{U_ERR}, and \code{V_ERR}.
Each archive carries a JSON string \code{meta_json} with the family, the sizes, and the constraint mode.

Load an archive from a short Python session:

\begin{lstlisting}
import json
import numpy as np

path = "analysis/optuna_manual_classical_dambreak/data/dam_break_elevation_compare_data.npz"
with np.load(path, allow_pickle=False) as z:
    model = z["M"]
    reference = z["A"]
    meta = json.loads(str(z["meta_json"]))
print(meta["model_family"], model.shape)
print(float(np.max(np.abs(model - reference))))
\end{lstlisting}

The trainer does not write PDF figures.
\file{swe_pinn/plots/jcp.py} contains helpers for loss curves and field panels.
The report script in the next section is the path that turns a finished study into figures.

\section{The study report}

From the repository root, with the database named in \file{user_input.py}:

\begin{lstlisting}[style=sh]
python analysis/analyze_study.py
\end{lstlisting}

The script reads that study, writes tables and figures under \file{analysis/<database name>/}, and compares the lowest-residual completed trial with the exact solution.
Override the locations when the active input file points somewhere else:

\begin{lstlisting}[style=sh]
python analysis/analyze_study.py \
    --db optuna_manual_eddies_three_family.db \
    --study manual_eddies_three_family \
    --problem steady_eddies
\end{lstlisting}

\code{--skip-train} writes the tradeoff report and the Optuna figures and skips the field comparison.
\code{--target} draws an optional horizontal line at that recorded objective in the text report.

\subsection{Text}

\file{data/param_tradeoff.txt} ranks completed trials by the recorded residual and lists family, size, and parameter counts.
That residual is the validation loss.
It is not the analytic field error.

\file{data/analytic_comparison.txt} is written when the comparison step runs.
It names the trial, the family, the architecture, the parameter counts, whether the weights were loaded or trained, the mesh, and for each field the discrete \(L_{2}\) norm of the error, the relative \(L_{2}\) error, and the \(L_{\infty}\) error.
\file{data/analytic_comparison.npz} stores the fields used for those numbers.

If saved weights match the selected architecture, they are loaded and not trained again.
If they are missing or were written for a different family or size, the script trains that architecture for \code{final_adam_epochs} on \code{final_adam_nx} by \code{final_adam_nt}, then saves the new weights.
A fixed-architecture example has no study, so this script has nothing to read until a search has completed at least one trial.

\subsection{Figures}

The script writes the following PDFs in \file{figures/} when the corresponding data exist.

\begin{center}
\begin{tabular}{@{}lp{8.6cm}@{}}
\toprule
File & Contents \\
\midrule
\file{optuna_objective_history.pdf} &
Objective against trial number \\
\file{optuna_error_vs_parameters.pdf} &
Residual against trainable-parameter count \\
\file{optuna_pareto_front.pdf} &
The nondominated residual-versus-size front \\
\file{optuna_model_family_comparison.pdf} &
Families side by side \\
\file{optuna_qubits_vs_error.pdf} &
Best residual at each qubit count \\
\file{optuna_quantum_layers_vs_error.pdf} &
Best residual at each layer count \\
\file{optuna_parameter_importance.pdf} &
fANOVA importance inside each family, one panel per family, each panel normalized on its own \\
\file{model_vs_analytic.pdf} &
Model, reference, and error for each field \\
\bottomrule
\end{tabular}
\end{center}

Family colors and markers are fixed in the script: classical is a blue circle, hybrid quantum is an orange square, and full quantum is a green triangle.
A reupload series is drawn separately when a study contains both embeddings.
A study that fixes \code{reupload}, as the shipped files do, is reported as the reupload series of that family.
Importance is skipped for a family with too few completed trials for fANOVA, and the script prints that it skipped the panel.

On \file{model_vs_analytic.pdf} each field is a row.
The columns are the model, the reference, and the signed error.
The horizontal axis is time for a two-field problem and \(y\) for a three-field problem.
The vertical axis is \(x\).
The error panel is annotated with the relative \(L_{2}\) error of that field.

\section{How to read a three-family figure}

On the error-versus-parameters figure, a family that sits lower at the same horizontal position reached a smaller residual at that size.
The parameter window in Section~\ref{sec:window} is what makes that horizontal comparison possible: draws outside the overlap were discarded, so the plotted points lie in a range every family was allowed to enter.
Full-quantum points still tend to lie toward the left of that window, because a pure circuit with a handful of qubits is smaller than a width-5, depth-7 perceptron.
Classical points can occupy the whole window.
Hybrid points occupy the part of the window that remains after the circuit is added to the trunk.

The Pareto figure is the set of completed trials that are not worse in both residual and size than some other completed trial.
A family can be absent from that front.
That is a result about the study you ran, at the epoch budget you set.
A 50-epoch smoke test does not support the same reading as a 4000-epoch study.

The analytic page is one model, the lowest residual in the study, retrained or reloaded at the final-stage budget.
Quote the trial number and the family from \file{analytic_comparison.txt} together with the relative \(L_{2}\) numbers.
The Optuna objective of that trial is a residual on the score lattice.
The relative \(L_{2}\) number is a field error on the evaluation mesh.
They answer different questions.
